\section{Introduction}
\IEEEPARstart{B}{etween-participant} variability is the main obstacle to EEG decoders that generalize to people absent from training \cite{saha2020variability,wu2022transfer}. In affective ERP data the participant factor carries several times the variance of the experimental condition. A family of label-free alignment methods addresses this by transforming each participant's data with statistics of that participant alone. Euclidean alignment whitens each participant's trials by their mean spatial covariance \cite{he2020ea}, and Riemannian methods recenter covariance matrices on the manifold \cite{zanini2018transfer,rodrigues2019rpa}. Subject centering is the first-order member of the family. It subtracts the participant's mean feature vector, and in repeated-measures ERP designs it produces large accuracy gains.

All of these methods are transductive. They use unlabeled observations of the test participant, and a reported gain describes deployment only if those observations exist when the decoder is applied. Three questions decide what a centering gain means, and the literature leaves them open. How many reference observations does the estimate need? When does centering lower accuracy? And what does a complete within-participant panel contribute beyond the offset estimate, given that its composition is fixed by the experimental design?

This paper answers them with a single additive model. The contributions are a reference-size law for the separation after centering (Proposition~\ref{prop:law}), a break-even reference size below which centering lowers accuracy (Corollary~\ref{cor:breakeven}), exact identities for the complete balanced panel, including the equivalence of leave-one-out and full-panel centering (Proposition~\ref{prop:panel}), and tests of each prediction on five public ERP CORE paradigms \cite{kappenman2021erpcore} and on trial-averaged affective ERPs of 211 participants, with Euclidean and Riemannian alignment evaluated on the same reference sets. The work is the geometric-transparency level of the ARSPI-Net framework \cite{lane2026dissertation,lane2026bibm} stated as a general result about EEG features.
